\documentclass[10pt]{extarticle}
\usepackage[a4paper,margin=1.5cm]{geometry}
\usepackage[T1]{fontenc}
\usepackage{amsmath,amssymb,mathtools}
\usepackage{graphicx}
\usepackage[export]{adjustbox}
\usepackage{multicol}
\usepackage{enumitem}
\usepackage{booktabs,array,tabularx}
\setlength{\columnsep}{0.7cm}
\setlength{\parindent}{0pt}
\setlength{\parskip}{3pt}
\pagestyle{empty}
\makeatletter
\newcommand{\qfigdef}[4]{\@namedef{qf@file@#1}{#2}\@namedef{qf@w@#1}{#3}\@namedef{qf@wide@#1}{#4}}
\newcommand{\qf@check}[1]{\@ifundefined{qf@file@#1}{\PackageError{score2}{Unknown figure id #1}{}}{}}
\newcommand{\qf@img}[4][]{\includegraphics[width=#2\dimexpr\@nameuse{qf@w@#4}pt\relax,max width=#3,max height=0.45\textheight,#1]{\@nameuse{qf@file@#4}}}
\newif\ifqf@full
\newif\ifqf@boxed
\newsavebox{\qf@box}
\newcommand{\qf@open}{\global\qf@fulltrue\par\global\qf@boxedtrue
  \setbox\qf@box\vbox\bgroup\hsize\textwidth\linewidth\textwidth}
\newcommand{\qf@emit}[1]{\par\egroup\global\qf@boxedfalse
  \ifdim\dimexpr\ht\qf@box+\dp\qf@box\relax<#1\textheight\noindent\usebox{\qf@box}\par\else\unvbox\qf@box\fi}
\newcommand{\qfullsplit}{\ifqf@boxed\qf@emit{0.85}\fi}
\newcommand{\qf@close}{\ifqf@boxed\qf@emit{0.45}\fi\global\qf@fullfalse}
\newcommand{\qfullbegin}{\end{multicols}\qf@open}
\newcommand{\qfullend}{\qf@close\begin{multicols}{2}}
\newcommand{\qfullstart}{\qf@open}
\newcommand{\qfulldone}{\qf@close}
\newcommand{\qfig}[2][1]{\qf@check{#2}%
  \ifqf@full\qf@blk{#1}{#2}\else
  \ifnum\@nameuse{qf@wide@#2}=1 \expandafter\@firstoftwo\else\expandafter\@secondoftwo\fi
  {\end{multicols}\begin{center}\qf@img{#1}{\textwidth}{#2}\end{center}\begin{multicols}{2}}%
  {\qf@blk{#1}{#2}}\fi}
\newcommand{\qf@blk}[2]{\par\smallskip{\centering\qf@img{#1}{\linewidth}{#2}\par}\smallskip}
\newcommand{\qfigi}[2][1]{\qf@check{#2}\qf@img[valign=c]{#1}{\linewidth}{#2}}
\makeatother
\newcommand{\papersec}[1]{\par\bigskip{\centering\Large\bfseries PAPER-#1\par}\medskip}
\newcommand{\partsec}[2]{\par\medskip\noindent{\large\bfseries PART-#1: #2}\par\smallskip}
\newcommand{\sect}[1]{\par\smallskip\noindent\textbf{#1}\par}
\newcommand{\qn}[1]{\par\smallskip\noindent\textbf{#1.}~}
\newcommand{\paraq}[2]{\par\smallskip\noindent\textbf{Paragraph for Questions #1 and #2}\par}
\newenvironment{qoptions}{\begin{itemize}[leftmargin=1.8em,itemsep=1pt,topsep=2pt,label={}]}{\end{itemize}}
\newcommand{\opt}[2]{\item[(#1)] #2}
\newcolumntype{L}[1]{>{\raggedright\arraybackslash}p{#1}}
\newcommand{\fwstart}{\end{multicols}\par\noindent\begin{minipage}{\textwidth}\small\setlength{\tabcolsep}{4pt}\renewcommand{\arraystretch}{1.5}}
\newcommand{\fwend}{\end{minipage}\par\medskip\begin{multicols}{2}}
% figure registry (generated by score2_tools.py finalize)
\qfigdef{p1_phy_s1_q01_1}{SRG (Enthusiast) - Test 02 (04 Oct 2020)-img/p1_phy_s1_q01_1.png}{151.5}{0}
\qfigdef{p1_phy_s1_q02_1}{SRG (Enthusiast) - Test 02 (04 Oct 2020)-img/p1_phy_s1_q02_1.png}{137.3}{0}
\qfigdef{p1_phy_s1_q03_1}{SRG (Enthusiast) - Test 02 (04 Oct 2020)-img/p1_phy_s1_q03_1.png}{105.6}{0}
\qfigdef{p1_phy_s1_q04_1}{SRG (Enthusiast) - Test 02 (04 Oct 2020)-img/p1_phy_s1_q04_1.png}{132.6}{0}
\qfigdef{p1_phy_s1_q05_1}{SRG (Enthusiast) - Test 02 (04 Oct 2020)-img/p1_phy_s1_q05_1.png}{155.4}{0}
\qfigdef{p1_phy_s1_q07_1}{SRG (Enthusiast) - Test 02 (04 Oct 2020)-img/p1_phy_s1_q07_1.png}{66.7}{0}
\qfigdef{p1_phy_s1_q09_1}{SRG (Enthusiast) - Test 02 (04 Oct 2020)-img/p1_phy_s1_q09_1.png}{115.2}{0}
\qfigdef{p1_phy_s1_q09_A}{SRG (Enthusiast) - Test 02 (04 Oct 2020)-img/p1_phy_s1_q09_A.png}{114.4}{0}
\qfigdef{p1_phy_s1_q09_B}{SRG (Enthusiast) - Test 02 (04 Oct 2020)-img/p1_phy_s1_q09_B.png}{114.4}{0}
\qfigdef{p1_phy_s1_q09_C}{SRG (Enthusiast) - Test 02 (04 Oct 2020)-img/p1_phy_s1_q09_C.png}{114.4}{0}
\qfigdef{p1_phy_s1_q09_D}{SRG (Enthusiast) - Test 02 (04 Oct 2020)-img/p1_phy_s1_q09_D.png}{114.4}{0}
\qfigdef{p1_phy_s1_q10_1}{SRG (Enthusiast) - Test 02 (04 Oct 2020)-img/p1_phy_s1_q10_1.png}{78.1}{0}
\qfigdef{p1_phy_s2_q12_1}{SRG (Enthusiast) - Test 02 (04 Oct 2020)-img/p1_phy_s2_q12_1.png}{138.7}{0}
\qfigdef{p1_phy_s2_q13_1}{SRG (Enthusiast) - Test 02 (04 Oct 2020)-img/p1_phy_s2_q13_1.png}{79.6}{0}
\qfigdef{p1_phy_s2_q14_1}{SRG (Enthusiast) - Test 02 (04 Oct 2020)-img/p1_phy_s2_q14_1.png}{124.4}{0}
\qfigdef{p1_phy_s2_q15_1}{SRG (Enthusiast) - Test 02 (04 Oct 2020)-img/p1_phy_s2_q15_1.png}{87.7}{0}
\qfigdef{p1_phy_s2_q16_1}{SRG (Enthusiast) - Test 02 (04 Oct 2020)-img/p1_phy_s2_q16_1.png}{104.7}{0}
\qfigdef{p1_phy_s2_q17_1}{SRG (Enthusiast) - Test 02 (04 Oct 2020)-img/p1_phy_s2_q17_1.png}{206.0}{0}
\qfigdef{p1_chem_s1_q05_1}{SRG (Enthusiast) - Test 02 (04 Oct 2020)-img/p1_chem_s1_q05_1.png}{281.5}{1}
\qfigdef{p1_chem_s1_q08_1}{SRG (Enthusiast) - Test 02 (04 Oct 2020)-img/p1_chem_s1_q08_1.png}{274.7}{1}
\qfigdef{p1_chem_s1_q08_2}{SRG (Enthusiast) - Test 02 (04 Oct 2020)-img/p1_chem_s1_q08_2.png}{124.2}{0}
\qfigdef{p1_chem_s1_q09_A}{SRG (Enthusiast) - Test 02 (04 Oct 2020)-img/p1_chem_s1_q09_A.png}{63.7}{0}
\qfigdef{p1_chem_s1_q09_B}{SRG (Enthusiast) - Test 02 (04 Oct 2020)-img/p1_chem_s1_q09_B.png}{64.4}{0}
\qfigdef{p1_chem_s1_q09_C}{SRG (Enthusiast) - Test 02 (04 Oct 2020)-img/p1_chem_s1_q09_C.png}{62.6}{0}
\qfigdef{p1_chem_s1_q09_D}{SRG (Enthusiast) - Test 02 (04 Oct 2020)-img/p1_chem_s1_q09_D.png}{56.1}{0}
\qfigdef{p1_chem_s1_q10_1}{SRG (Enthusiast) - Test 02 (04 Oct 2020)-img/p1_chem_s1_q10_1.png}{282.5}{1}
\qfigdef{p1_chem_s1_q10_3}{SRG (Enthusiast) - Test 02 (04 Oct 2020)-img/p1_chem_s1_q10_3.png}{399.9}{1}
\qfigdef{p1_chem_s2_q12_1}{SRG (Enthusiast) - Test 02 (04 Oct 2020)-img/p1_chem_s2_q12_1.png}{216.5}{0}
\qfigdef{p1_chem_s2_q14_1}{SRG (Enthusiast) - Test 02 (04 Oct 2020)-img/p1_chem_s2_q14_1.png}{416.3}{1}
\qfigdef{p1_chem_s2_q15_1}{SRG (Enthusiast) - Test 02 (04 Oct 2020)-img/p1_chem_s2_q15_1.png}{245.2}{0}
\qfigdef{p1_chem_s2_q15_11}{SRG (Enthusiast) - Test 02 (04 Oct 2020)-img/p1_chem_s2_q15_11.png}{195.9}{0}
\qfigdef{p1_chem_s2_q16_1}{SRG (Enthusiast) - Test 02 (04 Oct 2020)-img/p1_chem_s2_q16_1.png}{210.7}{0}
\qfigdef{p1_chem_s2_q18_1}{SRG (Enthusiast) - Test 02 (04 Oct 2020)-img/p1_chem_s2_q18_1.png}{195.1}{0}
\raggedcolumns\makeatletter\@beginparpenalty=10000\makeatother

\begin{document}
\begin{center}{\Large\bfseries SRG (Enthusiast) --- Test 02 (04 Oct 2020)}\end{center}
\begin{multicols}{2}
\partsec{1}{PHYSICS}
\sect{SECTION-I}
\qn{1} A particle of mass $m$ and charge $q$ moving with velocity $v$ enters Region II normal to the boundary as shown in the figure. Region II having a uniform magnetic field $B$ is perpendicular to the plane of the paper. The length of the Region II is $\ell$. Choose the correct choices.
\qfig{p1_phy_s1_q01_1}
\begin{qoptions}
\opt{A}{The particle enters Region III only if its velocity $v>\dfrac{q\ell B}{m}$}
\opt{B}{The particle enters Region III only if its velocity $v<\dfrac{q\ell B}{m}$}
\opt{C}{Path length of the particle in Region II is maximum when velocity $v<\dfrac{q\ell B}{m}$}
\opt{D}{Time spent in Region II is same for any velocity $v$ as long as the particle returns to Region I}
\end{qoptions}
\qn{2} Two long thin, parallel conductors carrying equal currents in the same direction are fixed parallel to the $x$-axis, one passing through $y=a$ and the other through $y=-a$. The resultant magnetic field due to the two conductors at any point is $B$. Which of the following are correct?
\qfig{p1_phy_s1_q02_1}
\begin{qoptions}
\opt{A}{$B=0$, for all points on $X$-axis}
\opt{B}{At all points on the $Y$-axis, excluding the origin, $B$ has only a $z$-component}
\opt{C}{At all points on $Z$-axis, excluding the origin, $B$ has only a $y$-component}
\opt{D}{$B$ can't have an $x$-component}
\end{qoptions}
\qn{3} From a cylinder of radius $R$, a cylinder of radius $R/2$ is removed as shown in the figure. Current flowing in the remaining cylinder is $I$. Magnetic field strength is
\qfig{p1_phy_s1_q03_1}
\begin{qoptions}
\opt{A}{Zero at point $A$}
\opt{B}{Zero at point $B$}
\opt{C}{$\dfrac{\mu_0I}{3\pi R}$ at point $A$}
\opt{D}{$\dfrac{\mu_0I}{3\pi R}$ at point $B$}
\end{qoptions}
\qn{4} A wire $B$ of finite length $L$ is kept on the right hand side of a long wire $A$ as shown in the figure. Direction of currents on both the wires are shown in figure. Suppose $F$ is the force on wire $B$ and $\tau$ the torque on it. Then,
\qfig{p1_phy_s1_q04_1}
\begin{qoptions}
\opt{A}{$F$ is upward, $F=\dfrac{\mu_0I_1I_2}{2\pi}\ln\!\left(\dfrac{L+d}{d}\right)$}
\opt{B}{$F$ is downward, $F=\dfrac{\mu_0I_1I_2}{2\pi}\ln\!\left(\dfrac{L+d}{d}\right)$}
\opt{C}{$\tau$ is clockwise}
\opt{D}{$\tau$ is anti-clockwise}
\end{qoptions}
\qn{5} The rectangular coil having 100 turns is placed in a uniform magnetic field of $\dfrac{0.05}{\sqrt2}\hat j$ as shown in the figure. Then,
\qfig{p1_phy_s1_q05_1}
\begin{qoptions}
\opt{A}{The magnetic dipole moment of the current carrying coil is $1.6\times10^{-2}\,\hat i\ (\text{A--m}^2)$}
\opt{B}{The torque acting on the coil is $5.66\times10^{-3}\ (\text{N--m})\,\hat k$}
\opt{C}{The torque acting on the coil is $0.566\times10^{-5}\ (\text{N--m})\,\hat k$}
\opt{D}{The magnetic dipole moment of coil is $16\times10^{-2}\,\hat i\ (\text{A--m}^2)$}
\end{qoptions}
\qn{6} A current carrying ring with its centre at origin and moment of inertia $2\times10^{-2}$ kg--m$^2$ about an axis passing through the centre and perpendicular to the plane has magnetic moment $M=(3\hat i-4\hat j)$ A--m$^2$. At time $t=0$ a magnetic field $B=(4\hat i+3\hat j)$ T is switched on. Then,
\begin{qoptions}
\opt{A}{Angular acceleration of the ring at time $t=0$ is 2500 rad s$^{-2}$}
\opt{B}{Maximum angular velocity of the ring is $50\sqrt2$ rad s$^{-1}$}
\opt{C}{Angular acceleration of the ring at time $t=0$ is 5000 rad s$^{-2}$}
\opt{D}{Maximum angular velocity of the ring is $25\sqrt2$ rad s$^{-1}$}
\end{qoptions}
\qn{7} A charged particle $A$ of charge $q=2$ C has velocity $v=100$ m/s. When it passes through point $A$ and has velocity in the direction shown in the figure. The strength of magnetic field at point $B$ due to this moving charge is ($r=2$ m)
\qfig{p1_phy_s1_q07_1}
\begin{qoptions}
\opt{A}{2.5 $\mu$T}
\opt{B}{5.0 $\mu$T}
\opt{C}{2.0 $\mu$T}
\opt{D}{None of these}
\end{qoptions}
\qn{8} A particle of specific charge $\alpha$ starts moving from the origin under the action of an electric field $E=E_0\hat i$ and magnetic field $B=B_0\hat k$. Its velocity at $(x_0,y_0,0)$ is $(4\hat i+3\hat j)$. The value of $x_0$ is
\begin{qoptions}
\opt{A}{$\dfrac{13}{2}\cdot\dfrac{\alpha E_0}{B_0}$}
\opt{B}{$\dfrac{16\alpha B_0}{E_0}$}
\opt{C}{$\dfrac{25}{2\alpha E_0}$}
\opt{D}{$\dfrac{5\alpha}{2B_0}$}
\end{qoptions}
\qn{9} A magnetic field $B=B_0\hat j$ exists in the region $a<x<2a$ and $B=-B_0\hat j$ in the region $2a<x<3a$, where $B_0$ is a positive constant. A positive point charge moving with a velocity $v=v_0\hat i$, where $v_0$ is a positive constant enters the magnetic field at $x=a$. The trajectory of the charge in this region can be like
\qfig{p1_phy_s1_q09_1}
\begin{qoptions}
\opt{A}{\qfigi{p1_phy_s1_q09_A}}
\opt{B}{\qfigi{p1_phy_s1_q09_B}}
\opt{C}{\qfigi{p1_phy_s1_q09_C}}
\opt{D}{\qfigi{p1_phy_s1_q09_D}}
\end{qoptions}
\qn{10} Find the magnetic field at point $P$ due to arrangement as shown in the figure.
\qfig{p1_phy_s1_q10_1}
\begin{qoptions}
\opt{A}{$\dfrac{\mu_0I}{\sqrt2\pi d}\left(1-\dfrac1{\sqrt2}\right)$}
\opt{B}{$\dfrac{2\mu_0I}{\sqrt2\pi d}$}
\opt{C}{$\dfrac{\mu_0I}{\sqrt2\pi d}$}
\opt{D}{$\dfrac{\mu_0I}{\sqrt2\pi d}\left(1+\dfrac1{\sqrt2}\right)$}
\end{qoptions}
\sect{SECTION-II}
\qn{11} Two parallel wires in the plane of the paper are distance $X_0$ apart. A point charge is moving with speed $u$ between the wires in the same plane at a distance $X_1$ from one of the wires. When the wires carry current of magnitude $I$ in the same direction, the radius of curvature of the path of the point charge is $R_1$. In contrast, if the current $I$ in the two wires have directions opposite to each other, the radius of curvature of the path is $R_2$. If $\dfrac{X_0}{X_1}=3$, then find the value of $\dfrac{R_1}{R_2}$
\qn{12} As shown in the figure, three sided frame is pivoted at $P$ and $Q$ and hangs vertically. Its sides are of same length and have a linear density of $\sqrt3$ kg/m. A current of $10\sqrt3$ A is sent through the frame, which is in a uniform magnetic field of 2 T directed upwards as shown in the figure. Then, angle through which the frame will be deflected in equilibrium is $\dfrac{\pi}{k}$, Find the value of $k$. (Given, $g=10$ m/s$^2$)
\qfig{p1_phy_s2_q12_1}
\qn{13} A non-uniform magnetic field $B=B_0\left(1+\dfrac yd\right)\hat k$ is present in region of space in between $y=0$ and $y=d$. The lines are shown in the figure. A particle of mass $m$ and positive charge $q$ is moving. Given, an initial velocity $v=v_0\hat i$. The $x$-component of velocity of the particle is $V_x=V_0-\dfrac{kqB_0d}{2m}$ when it leaves the field, then find the value of $k$.
\qfig{p1_phy_s2_q13_1}
\qn{14} For the given circuit shown,
\qfig{p1_phy_s2_q14_1}
If a voltmeter with sensitivity of $1\,\mathrm{k\Omega/V}$ is connected between points $A$ and $B$, then percentage error in the reading is found $10\times k$. Find the value of $k$.
\qn{15} A non-conducting sphere has mass of 100 g and radius 20 cm. A flat compact coil of wire with 5 turns is wrapped tightly around it with each turns concentric with the sphere. This sphere is placed on an inclined plane such that plane of coil is parallel to the inclined plane. A uniform magnetic field of 0.5 T exists in the region in vertically upwards direction. If a current of $\dfrac k\pi$ A is required to rest the sphere in equilibrium, then find the value of $k$.
\qfig{p1_phy_s2_q15_1}
\qn{16} A current carrying small circular loop is placed at a distance 2 cm from an infinitely long current carrying conductor as shown in the figure. Wire and loop lies in same plane. The loop has 100 turns, its effective radius is 2 cm. The torque acting on the loop in the given situation is $x\mu_0$. Find the value of $x$.
\qfig{p1_phy_s2_q16_1}
\qn{17} A stiff metal rod is kept over two knife edges of length $L=1$ m as shown in the figure.
\qfig{p1_phy_s2_q17_1}
The rod carries a current of 16 A (in the direction shown) and rolls along the rails without slipping due to the magnetic force produced by a uniform magnetic field of 0.5 T perpendicular and pointing downwards. If rod starts from rest and its speed when it leaves the rails is $\dfrac{k}{\sqrt5}$ ms$^{-1}$, then find the value of $k$.
\qn{18} Let a wire is placed in a magnetic field that varies with distance from origin as $B=B_0\left(1+\dfrac xa\right)\hat k$. Ends of wire are at $(a,0)$ and $(2a,0)$ and its carries a current $i$. If force on wire is $F=B_0i\left(\dfrac{ka}{2}\right)(-\hat j)$. Then, find the value of $k$.
\qn{19} A long conducting cylinder of radius $R$ carries a current $I$. The current density $J$ is a function of radius $r$ as $J=br$, where $b$ is a constant. Magnetic field at a distance $r>R$ measured from the axis is $B=\dfrac{\mu_0bR^3}{Nr}$, then find the value of $N$.
\qn{20} A packed bundle of 100 long straight, insulated wires forms a cylinder of radius $R=0.5$ cm. If each wire carries a current of 2 A, then magnetic force per unit length on a wire located at 2 mm from the centre of the bundle is $N\times10^{-4}$ N/m. Find the value of $N$.
\partsec{2}{CHEMISTRY}
\sect{SECTION-I}
\qn{1} The incorrect statement is
\begin{qoptions}
\opt{A}{Vapour pressure of a liquid always increases by increasing temperature}
\opt{B}{Vapour pressure only depends upon temperature and not on the nature of substance}
\opt{C}{Vapour pressure doesnot depends upon the quality of liquid taken and the surface area of the liquid}
\opt{D}{Vapour pressure is not a colligative property and is independent of concentration of liquid}
\end{qoptions}
\qn{2} Relative decrease in vapour pressure of an aq solution containing 2 moles $\left[\mathrm{Cu(NH_3)_2Cl}\right]\mathrm{Cl}$ in 3 moles of $\mathrm{H_2O}$ is 0.50. On reaction with $\mathrm{AgNO_3}$, this solution will form (Assuming no change in degree of ionisation of substance on adding $\mathrm{AgNO_3}$)
\begin{qoptions}
\opt{A}{1 mol AgCl}
\opt{B}{0.25 mol AgCl}
\opt{C}{0.5 mol AgCl}
\opt{D}{0.40 mol AgCl}
\end{qoptions}
\qn{3} An ideal mixture of liquids with 2 moles A and 2 moles of B has total vapour pressure 1 atm at constant temperature. Another mixture with 1 mol A and 3 mol B has a vapour pressure greater than 1 atm. But if 4 moles of C are added to second mixture the vapour pressure comes down to 1 atm. Vapour pressure of pure A and pure B will be :

[Vapour pressure of C, $P_C^0=0.8$ atm]
\begin{qoptions}
\opt{A}{$P_A^0=1.4$ atm, $P_B^0=0.7$ atm}
\opt{B}{$P_A^0=1.2$ atm, $P_B^0=0.6$ atm}
\opt{C}{$P_A^0=1.4$ atm, $P_B^0=0.6$ atm}
\opt{D}{$P_A^0=0.6$ atm, $P_B^0=1.4$ atm}
\end{qoptions}
\qn{4} A solute `S' undergoes reversible trimerisation when dissolved in a certain solvent. The boiling point elevation of its 0.1 molal solution was found to be identical to the boiling point elevation in case of a 0.08 molal solution of a solute which neither dissociate nor associate. What will be the percenatge trimerisation of solute `S'
\begin{qoptions}
\opt{A}{30\%}
\opt{B}{40\%}
\opt{C}{50\%}
\opt{D}{60\%}
\end{qoptions}
\qfullbegin
\qn{5} Dry air is slowly passed through three solutions of different concentrations $C_1, C_2$ and $C_3$, each containing nonvolatile NaCl as solute and water as solvent as shown in figure. If vessel 2 gains weight and vessel 3 loses weight then
\qfig{p1_chem_s1_q05_1}
\qfullsplit
\begin{qoptions}
\opt{A}{$C_1>C_2$}
\opt{B}{$C_1<C_2$}
\opt{C}{$C_3>C_1$}
\opt{D}{$C_2>C_3$}
\end{qoptions}
\qfullend
\qn{6} In which of the following pairs of solutions the value of Vann-Hoff factor will be same (considering 100\% dissociation)
\begin{qoptions}
\opt{A}{0.05 M $\mathrm{K_4[Fe(CN)_6]}$ and 0.1M $\mathrm{FeSO_4}$}
\opt{B}{0.1 M $\mathrm{K_4[Fe(CN)_6]}$ and\\ 0.05 M FeSO$_4$(NH$_4$)$_2$SO$_4{\cdot}$6H$_2$O}
\opt{C}{0.2 M NaCl and 0.1 M $\mathrm{BaCl_2}$}
\opt{D}{0.05 M $\mathrm{FeSO_4(NH_4)_2SO_4{\cdot}6H_2O}$ and 0.02 M $\mathrm{KCl{\cdot}MgCl_2{\cdot}6H_2O}$}
\end{qoptions}
\qn{7} Two liquids A and B form an ideal solution. The vapour pressure of solution is 700 torr at 80$^\circ$C. It is distilled till $\dfrac{2^{rd}}{3}$ of the solution is collected as condensate. The composition of the condensate is $X_A'=0.75$ and that of residue is $X_A''=0.30$. If the vapour pressure of the residue is 600 torr at 80$^\circ$C, then which of the following is / are true
\begin{qoptions}
\opt{A}{The composition of original mixture was $X_A=0.6$}
\opt{B}{The composition of original mixture was $X_A=0.4$}
\opt{C}{$P_A^0=\dfrac{2500}{3}$ torr}
\opt{D}{$P_B^0=500$ torr}
\end{qoptions}
\qfullbegin
\qn{8} SPM is the membrane which permits the solvent to move across it. An experiment is done with 0.1 M 500 ml aq solution of $\mathrm{CH_3COOH}$ and 0.1 M 750 ml aq solution of NaOH as given figure
\qfig{p1_chem_s1_q08_1}
\qfullsplit

Neutralisation reaction is $\mathrm{CH_3COOH}+\mathrm{NaOH}\to\mathrm{CH_3COONa}+\mathrm{H_2O}$

Choose the incorrect statement / s
\begin{qoptions}
\opt{A}{Neutralisation occurs in container A}
\opt{B}{Neutralisation occurs in container B}
\opt{C}{Meq of salt $\mathrm{CH_3COONa}$ formed will be 50}
\opt{D}{Solution behaves like a buffer}
\end{qoptions}
\qfullend
\qn{9} \qfig{p1_chem_s1_q08_2}
Out of the following which one / is are correct
\begin{qoptions}
\opt{A}{1, \qfigi{p1_chem_s1_q09_A}}
\opt{B}{2, \qfigi{p1_chem_s1_q09_B}}
\opt{C}{3, \qfigi{p1_chem_s1_q09_C}}
\opt{D}{4, \qfigi{p1_chem_s1_q09_D}}
\end{qoptions}
\qfullbegin
\qn{10} \qfig{p1_chem_s1_q10_1}
\qfullsplit
$X, Y, Z$ are :
\qfig{p1_chem_s1_q10_3}
\qfullend
\sect{SECTION-II}
\qn{11} 25 ml of $\mathrm{FeC_2O_4}$ is dissolved in 186 gm of water. Calculate the depression in freezing point if 10 ml of same $\mathrm{FeC_2O_4}$ are titrated exactly with 30 ml of 0.4 M acidic $\mathrm{KMnO_4}$ solution.

[$K_f\,\mathrm{H_2O}=1.86$ unit assume 100\% ionisation of $\mathrm{FeC_2O_4}$]
\qn{12} CD : Fixed wall EF : Semi permeable memberance, fixed wall allows the pressure transfer but no mass transfer. If 16 gm He present in ABCD chamber. Then calculate the length of this chamber in cm if no osmosis takes place through SPM `EF'. Assume 75\% ionisation $\mathrm{K_4\left[Fe(CN)_6\right]}$
\qfig{p1_chem_s2_q12_1}
\qn{13} The vapour pressure of Fluorobenzene at $t^\circ$C is given by the equation $\left[\log p(\mathrm{mmHg})=7.0-\dfrac{1250}{t+220}\right]$ then calculate boiling point of liquid in $^\circ$C if the external applied pressure is 5.26 \% more then required for normal boiling point [$\log 2=0.3$]
\qfullbegin
\qn{14} Find the number of compounds which on reaction with $\mathrm{LiAlH_4}$ followed by treatment with water give optically inactive compound ?
\qfig{p1_chem_s2_q14_1}
\qfullend
\qn{15} In how many of the following reactions, reagents are wrongly chosen ?
\qfig[0.75]{p1_chem_s2_q15_1}
\qfig{p1_chem_s2_q15_11}
\qn{16} A decapeptide of molecular weight 796 g/mol on hydrolysis gives glycine, alanine and phenyl alanine, of which glycine contributes 47\% of the hydrolysed products by weight. What is the number of alanines in the decapeptide ?
\columnbreak
\qn{17} Number of cyclic product(s) formed
\qfig{p1_chem_s2_q16_1}
\qn{18}
\qfig{p1_chem_s2_q18_1}
(Assuming all the substrate convert into substitution products containing 0.05 mole of S-configuration). Calculate the percentage of $\mathrm{S_N2}$ mechanism.
\qn{19} It required 0.7 g of a hydrocarbon (A) to react completely with $\mathrm{Br_2}$ (2.0 g) and form a non resolvable product. On treatment of (A) with HBr it yielded monobromo alkane (B). The same compound (B) was obtained when (A) was treated with HBr in presence of peroxide. Write down the structure formula of (A) and (B) and explain the reactions involved. Total dibromoderivative possible for (A) are :
\qn{20} How many of the following are classified as LINEAR polymers ?

LDPE, PVC, amylopectin, cellulose, melamine, nylon 6, isoprene and polystyrene
\partsec{3}{MATHEMATICS}
\sect{SECTION-I}
\qn{1} The differential equation of the curve such that the $y$-intercept of any tangent is equal to the corresponding subnormal is:
\begin{qoptions}
\opt{A}{A linear equation}
\opt{B}{A homogenous equation}
\opt{C}{An equation with separable variables}
\opt{D}{None of these}
\end{qoptions}
\qn{2} The equation of the curve passing through the point $\left(1,\dfrac\pi4\right)$ and having slope of tangent at any $(x,y)$ as $\dfrac yx-\cos^2\dfrac yx$ is equal to
\begin{qoptions}
\opt{A}{$x=e^{1-\tan y/x}$}
\opt{B}{$y=x\tan^{-1}(\log x/e)$}
\opt{C}{$y=x\tan^{-1}\log(e/x)$}
\opt{D}{$x=e^{1+\cot y/2}$}
\end{qoptions}
\qn{3} If $f(x)$ is a function such that $x\displaystyle\int_0^x(1-t)f(t)\,dt=\int_0^x tf(t)\,dt$ and $f(1)=1$ then find $f(x)$
\begin{qoptions}
\opt{A}{$e^{-\frac1x-3\log x-1}$}
\opt{B}{$e^{\frac{-1}{x}}-3\log x+1$}
\opt{C}{$\dfrac1{x^3}e^{-\frac1x-1}$}
\opt{D}{None of these}
\end{qoptions}
\qn{4} For which of the following values of $m$, is the area of the region bounded by the curve $y=x-x^2$ and the line $y=mx$ equals $9/2$?
\begin{qoptions}
\opt{A}{$-4$}
\opt{B}{$-2$}
\opt{C}{2}
\opt{D}{4}
\end{qoptions}
\qn{5} Let $C_1: y=ex\ln x$ and $C_2: y=\dfrac{\ln x}{ex}$ be two curves, then
\begin{qoptions}
\opt{A}{Area bounded by these curves above the $x$-axis is zero}
\opt{B}{Area bounded by these curves below the $x$-axis is $\dfrac{e^2-4}{3e}$}
\opt{C}{Area bounded by these curves below the $x$-axis is $\dfrac{e^2-5}{4e}$}
\opt{D}{None of these}
\end{qoptions}
\qn{6} The solution of the differential equation $f(x)\dfrac{dy}{dx}+f'(x)y=1$ is
\begin{qoptions}
\opt{A}{$x=y\,f(x)+c$}
\opt{B}{$f^{-1}(x)+c=0$}
\opt{C}{$y=\dfrac{x+c}{f(x)}$}
\opt{D}{None of these}
\end{qoptions}
\qn{7} For the differential equation $y=e^{dy/dx}$
\begin{qoptions}
\opt{A}{Order is not defined}
\opt{B}{Order is 1}
\opt{C}{Degree is undefined}
\opt{D}{Degree is 1}
\end{qoptions}
\qn{8} The differential equation representing the family of curves $y^2=2c\left(x+\sqrt c\right)$, where $c$ is a positive parameter, is of
\begin{qoptions}
\opt{A}{Order - 1}
\opt{B}{Order - 2}
\opt{C}{Degree - 3}
\opt{D}{Degree - 4}
\end{qoptions}
\qn{9} Solution of differential equation $\left(e^x+1\right)y\,dy=(y+1)e^x dx$ is
\begin{qoptions}
\opt{A}{$\left(e^x+1\right)(y+1)=ce^y$}
\opt{B}{$\left(e^x+1\right)\left|y+1\right|=ce^y$}
\opt{C}{$\left(e^x+1\right)(y+1)=\pm ce^y$}
\opt{D}{None of these}
\end{qoptions}
\qn{10} The area bounded by a curve, the axis of co-ordinates \& the ordinate of some point of the curve is equal to the length of the corresponding arc of the curve. If the curve passes through the point $P(0,1)$ then the equation of this curve can be :
\begin{qoptions}
\opt{A}{$y=\dfrac12\left(e^x-e^{-x}+2\right)$}
\opt{B}{$y=\dfrac12\left(e^x+e^{-x}\right)$}
\opt{C}{$y=1$}
\opt{D}{$y=\dfrac{2}{e^x+e^{-x}}$}
\end{qoptions}
\sect{SECTION-II}
\qn{11} Let $\displaystyle\int_0^{e^x}\left(t^2-1\right)f(t)\,dt=\dfrac1{15}\left[3e^{5x}-5e^{3x}+15\right]$. Then the area bounded by $y=f(x)$ and $y=4$ is $m/n$ where $m$ and $n$ are coprime positive integers then find $m-10n$
\qn{12} Tangents (parallel to angle bisectors of lines $lx+my+n=0$ and $ax+by+c=0$) are drawn to ellipse $\dfrac{x^2}{36}+\dfrac{y^2}{64}=1$, so as to generate a quadrilateral with vertices as points of intersection of these tangents and inscribing the ellipse. If the maximum area of the quadrilateral is $n$. Find sum of digits of $[n]$ (where $[.]$ represents greatest integer function)
\qn{13} The area of the figure bounded by $y^2=9x$ and $y=3x$ is $m/n$ units where $m$ and $n$ are coprime positive integers find $m+n$
\qn{14} Area bounded by the curve $y=|x|-2$ and $y=1-|x-1|$ is equal to
\qn{15} Area of the triangle formed by the $x$-axis and tangent and normal drawn to the curve $y=x^2$ at $(1,1)$ is $m/n$ where $m$ and $n$ are coprime positive integers then $m+n$ is equal to
\qn{16} Area enclosed by the curve $y=f(x)$ that is being defined parametrically as $x=\dfrac{1-t^2}{1+t^2}$, $y=\dfrac{2t}{1+t^2}$ is equal to $k\pi$ sq. units then $k$ is equal to
\qn{17} The area bounded by the normal at $(1,2)$ on the curve, $y^2=4x$, $x$-axis and the curve is given by $m/n$ sq. units where $m$ and $n$ are coprime integers find $m-n$
\qn{18} The area of the figure bounded by $y=\sin x$, $y=\cos x$ and $x$-axis in the first quadrant is equal to $2\left(\sqrt m-n\right)$ where $m$ and $n$ are positive integers then find $m+n$
\qn{19} The degree of the differential equation $\left(\dfrac{d^3y}{dx^3}\right)^{2/3}+4-3\dfrac{d^2y}{dx^2}+5\dfrac{dy}{dx}=0$ is
\qn{20} The order of the differential equation whose general solution is given by $y=(c_1+c_2)\cos(x+c_3)-c_4e^{x+c_5}$, where $c_1,c_2,c_3,c_4,c_5$ are arbitrary constants, is
\end{multicols}
\end{document}
